\scanpage{44}
\mspage{21}
respecte également la structure (95). Il serait bon d'en comprendre la raison a priori.

\deleted{Pour ceci, je note que l'action d'un autom. \deleted{de} $u$ de $P_M$ est connue, quand on la connaît sur l'ens des sommets $S_{P_M}$, puisque les \deleted{\unsure{$S$}} ens des arêtes et \deleted{\unsure{$F$}} \unsure{ens} des faces s'identifient à des ens. de parties (à deux \deleted{\unsure{ou}} et à trois éléments respectivement) de $S_{P_M}$. D'autre part, un autom. \inserted{\unsure{$u$|``}}qui respecte la structure (101) est connu quand on connaît son action sur $\pi_0(P_M)$, \uline{et} sur \uline{un} $S_{M,i}$. Quand $u$ respecte \inserted{aussi} la structure (97), l'action de $u$ sur $\pi_0(P_M)$ se décrit via $\delta(u)\in Z'_\Lambda$. Lorsque $u$ est de la forme $T_{\dot\lambda}$, $\dot\lambda\in Z'_\Lambda$, on a $\delta(u)=\dot\lambda^2$, i.e l'action de $T_{\dot\lambda}$ sur $\pi_0(P_M)$ est déterminée par la structure (97).}
\margin{\uline{NB} On a en vertu des structures (97) (101) $S_{P_M}\simeq I_M\times\pi_0(P_M)$, où $I_M$ est ``la val. commune'' des $S_{P_M,i}$, identifiés via les \unsure{tr} de (101).}
% slanted

Récapitulons le type d'objet structuré qui est pour nous $P_M$. \deleted{Il}\unsure{Il|On} peut \unsure{se|le} décrire par les données
\[
\left\{\begin{aligned}
&\deleted{I}S\ \text{un ens }\deleted{\unsure{\ne\varnothing|\text{fini}}}\inserted{\text{(de cardinal $\geqslant3$)}}\
\text{(la ``valeur commune'' des }\deleted{\unsure{S_{M,i}}}S_{P_M,i},\ \text{pour }i\in\pi_0(P_M)),\\
&\pi\ \text{un }Z'_\Lambda\text{-torseur}\qquad(Z'_\Lambda=\Lambda^*/\{\pm1_\Lambda\})\\
&\pi\xrightarrow{\Sigma}\unsure{\mathfrak P|\mathcal P}(\mathfrak P_3(\deleted{I}S))\quad\text{une application}
\end{aligned}\right.\tag*{(105)}
\]
\margin{en fait $\mathrm{card}\,S\geqslant4$}
% slanted
satisfaisant les conditions nécessaires suivantes\,:
\[
\left\{\begin{aligned}
&\text{\textcircled{a}}\quad\forall\,i\in\pi,\quad \Sigma(i)\in\mathfrak P(\mathfrak P_3(\unsure{I|S|\mathbb S}))\ \text{est une structure}\\
&\hspace{4em}\text{de 2-polyèdre triangulé }\inserted{\text{connexe}}\ \text{sur }\deleted{I}S,\ \deleted{\text{soit}}\\
&\deleted{\text{b) Les }u\in\mathfrak S_I\ \text{qui sont des autom.}}
\end{aligned}\right.\tag*{(106)}
\]
% archive sheet number: 43
